\subsection{Stage 4: MEAM Initialization and Library Merging}
\label{sec:stage4}

The DFT reference is mapped onto MEAM parameters through analytic relations.
The single-element mappings are direct,
\begin{align*}
	\texttt{esub} &\leftarrow E_\mathrm{coh}, &
	\texttt{alat} &\leftarrow a_0,
\end{align*}
while the Rose shape parameter follows from the universal equation of state
\citep{vinet1989universal}:
\begin{equation}
	\alpha = \sqrt{\frac{9\,B\,\Omega}{E_\mathrm{coh}}}, \qquad
	\Omega = \begin{cases}
		a_0^3/4              & \text{FCC}, \\
		a_0^3/2              & \text{BCC}, \\
		a_0^3 \sqrt{2}       & \text{HCP (approx.)}, \\
		a_0^3/8              & \text{diamond},
	\end{cases}
	\label{eq:alpha}
\end{equation}
where $B$ is converted from GPa to eV/\AA{}$^3$ before substitution. For
aluminium ($a_0=4.047$~\AA, $E_\mathrm{coh}=3.611$~eV, $B=77.5$~GPa, FCC),
Eq.~\eqref{eq:alpha} gives $\Omega=16.57$~\AA{}$^3$,
$B_{\text{eV/\AA{}}^3}=0.484$, and $\alpha\approx 4.47$.

The cross-interaction terms are seeded from the binary formation energies and
the elemental nearest-neighbour distances:
\begin{align}
	E_c(i,j)     &= \tfrac{1}{2}\bigl[\texttt{esub}_i + \texttt{esub}_j\bigr] + E_\mathrm{form}(i,j), \\
	r_e(i,j)     &= \tfrac{1}{2}\bigl[d_\mathrm{nn}(i) + d_\mathrm{nn}(j)\bigr], \\
	\alpha(i,j)  &= \tfrac{1}{2}\bigl[\alpha_i + \alpha_j\bigr].
\end{align}

\paragraph{Automatic library discovery and merging.}
Because the public MEAM libraries each cover a different subset of elements,
Stage~4 starts with a discovery and merge pass:
\begin{enumerate}
	\item enumerate every available MEAM library file;
	\item find the smallest set of files whose union covers the requested
	      element list;
	\item parse the single-element entries from each chosen file and
	      concatenate them, keeping only the first occurrence of any duplicated
	      species;
	\item fill any cross-interaction entry missing from the merged set with the
	      DFT-seeded value;
	\item write the unified library and parameter files to a DFT-initialized
	      output location.
\end{enumerate}
The default twelve-element selection produces one merged library covering Al,
Co, Cr, Cu, Fe, Mg, Mn, Mo, Ni, Si, Ti, and Zn. As an example of the merge
logic, an Al--Cu--Zn--Mg--Fe--Cr--Ti potential draws Al, Cu, Fe, Cr, and Ti
from the Fe--Mn--Ni--Ti--Cu--Cr--Co--Al library, Mg and Zn from the
Al--Mg--Zn library, and synthesises the missing cross-terms (Cr--Mg, Ti--Zn,
and so on) from the DFT formation energies. The whole pass takes about
$70$\,ms and is cached.

\paragraph{Backfilling incomplete elemental data.}
A few elements never reach a clean cubic $C_{11}/C_{12}$. The non-cubic
lattices of Si, Ti, Zn, and Mg skip the cubic elastic-constant computation
entirely, and the magnetic Fe and Mn rows come out mechanically unstable
(Section~\ref{sec:stage3}). For these a helper backfills any missing or
unphysical $C_{ij}$ from the recorded bulk modulus $B$ and a per-element
Poisson-ratio default, so Stage~4 always receives a complete and
self-consistent set.

The single-element parameters produced by applying Eq.~\eqref{eq:alpha} to the
DFT reference are listed in Table~\ref{tab:meam_init}. The Mn row is flagged
because its EOS scan did not converge for the simplified $\alpha$-Mn
approximation; its $\alpha$ is a placeholder and is overwritten by the donor
library entry during the merge. The other rows feed directly into the
\texttt{esub}, \texttt{alat}, and \texttt{alpha} entries of the
DFT-initialized parameter file.

% AUTO-GENERATED by src/stats/ -- do not edit by hand
\begin{table}[htbp]
\centering
\caption{MEAM single-element initialization values computed from the DFT reference of Stage~3 by Eq.~\eqref{eq:alpha}.}
\label{tab:meam_init}
\begin{tabular}{llrrrr}
\hline
 Element   & Lattice   &   $a_0$ (\AA) &   $E_\mathrm{coh}$ (eV) &   $\Omega$ (\AA$^3$) &   $\alpha$ \\
\hline
 Mg        & HCP       &         3.177 &                   1.483 &               45.357 &      9.063 \\
 Mo        & BCC       &         3.169 &                  10.353 &               15.911 &      4.924 \\
 Al        & FCC       &         4.047 &                   3.611 &               16.569 &      4.470 \\
 Cu        & FCC       &         3.639 &                   3.857 &               12.051 &      4.930 \\
 Fe        & BCC       &         2.806 &                   6.834 &               11.046 &      6.179 \\
 Zn        & HCP       &         2.766 &                   0.964 &               29.931 &     10.921 \\
 Ti        & HCP       &         2.915 &                   6.284 &               35.022 &      5.987 \\
 Cr        & BCC       &         2.848 &                   3.964 &               11.549 &      5.906 \\
 Mn        & BCC       &         2.803 &                   3.728 &               11.012 &      6.120 \\
 Si        & DIAMOND   &         5.471 &                   5.130 &               20.464 &      4.473 \\
\hline
\end{tabular}
\end{table}

